\chapter{Tổng quan bài toán và các công trình liên quan}
\label{chap:tong-quan}

\section{Đặt vấn đề}

Trong những năm gần đây, máy bay không người lái, hay drone, ngày càng được sử dụng rộng rãi trong các bài toán vận tải, giao hàng, giám sát và kiểm tra hạ tầng. So với các phương tiện mặt đất truyền thống, drone có lợi thế là có thể di chuyển trực tiếp trong không gian, không bị ràng buộc bởi mạng đường bộ, đồng thời có khả năng tiếp cận những khu vực khó đi lại hoặc nguy hiểm đối với con người. Vì vậy, drone đặc biệt phù hợp với các nhiệm vụ kiểm tra đường dây điện, đường ống, mạng lưới giao thông, kênh mương hoặc các tuyến hạ tầng dạng đường.

Trong các ứng dụng kiểm tra hạ tầng tuyến tính, đối tượng cần phục vụ thường không phải là các điểm rời rạc mà là các cạnh hoặc các đoạn tuyến liên tục. Do đó, lớp bài toán phù hợp để mô hình hóa các ứng dụng này là bài toán định tuyến trên cung, hay \textit{Arc Routing Problem} (ARP), thay vì các bài toán định tuyến trên nút như Traveling Salesman Problem hoặc Vehicle Routing Problem. Trong bài toán định tuyến trên cạnh với drone, drone cần bay dọc theo các đoạn tuyến yêu cầu để thực hiện nhiệm vụ phục vụ, chẳng hạn như ghi hình, kiểm tra hoặc thu thập dữ liệu cảm biến.

Tuy nhiên, việc sử dụng drone cũng tạo ra nhiều thách thức mới. Thứ nhất, drone có giới hạn về thời lượng pin nên mỗi chuyến bay chỉ có thể kéo dài tối đa trong một khoảng thời gian hoặc quãng đường nhất định. Thứ hai, đối với các mạng hạ tầng lớn, một drone đơn lẻ không thể hoàn thành toàn bộ nhiệm vụ. Thứ ba, drone thường cần phương tiện mặt đất hỗ trợ để vận chuyển đến gần khu vực cần kiểm tra, thay pin, cất cánh và hạ cánh. Vì vậy, bài toán thực tế không chỉ là tìm đường bay cho drone, mà còn phải đồng thời quyết định vị trí triển khai phương tiện mặt đất, phân công các đoạn tuyến cần phục vụ cho từng drone, và xây dựng nhiều chuyến bay khả thi cho mỗi drone.

Từ những yêu cầu trên, bài toán \textit{Min--Max Multi--Trip Drone Location Arc Routing Problem} (MM-MT-dLARP) được đề xuất nhằm mô hình hóa tình huống trong đó có một depot trung tâm, một tập các xe tải, mỗi xe tải mang theo một drone, và một tập các điểm tiềm năng để xe tải có thể đến triển khai drone. Mỗi drone có thể thực hiện nhiều chuyến bay từ điểm triển khai của nó, nhưng mỗi chuyến bay phải thỏa mãn giới hạn về thời lượng bay. Mục tiêu của bài toán là hoàn thành toàn bộ nhiệm vụ kiểm tra trong thời gian sớm nhất, tức là tối thiểu hóa thời gian hoàn thành lớn nhất trong số các xe tải và drone tương ứng \cite{corberan2025mmdlarp}.

\section{Mô tả khái quát bài toán MM-MT-dLARP}

Bài toán MM-MT-dLARP xét một tập các đoạn tuyến cần được phục vụ bởi drone. Các đoạn tuyến này thường là các đường cong hoặc các đoạn hạ tầng liên tục trong không gian. Có một depot ban đầu, nơi xuất phát của \(P\) xe tải. Mỗi xe tải mang theo một drone và có thể di chuyển từ depot đến một điểm triển khai thuộc tập \(\mathcal{D}\). Mỗi điểm triển khai có thể được hiểu là một vị trí vận hành, nơi xe tải dừng lại để drone cất cánh, hạ cánh và thay pin nếu cần.

Mỗi drone có thể thực hiện nhiều chuyến bay từ điểm triển khai đã chọn. Một chuyến bay của drone bắt đầu tại điểm triển khai, bay đến một hoặc nhiều đoạn tuyến cần phục vụ, thực hiện nhiệm vụ trên các đoạn tuyến đó, sau đó quay trở lại đúng điểm triển khai ban đầu. Do giới hạn pin, chi phí hoặc thời gian của mỗi chuyến bay không được vượt quá một hằng số \(L\). Ngoài ra, hai xe tải khác nhau không được sử dụng cùng một điểm triển khai.

Mục tiêu của bài toán là lựa chọn \(P\) điểm triển khai trong tập \(\mathcal{D}\), phân công các đoạn tuyến cần phục vụ cho các drone, và xây dựng tập các chuyến bay cho từng drone sao cho toàn bộ các đoạn tuyến yêu cầu được phục vụ. Hàm mục tiêu có dạng min--max: với mỗi xe tải, chi phí được tính bằng chi phí di chuyển từ depot đến điểm triển khai và quay lại, cộng với tổng chi phí các chuyến bay của drone tương ứng; bài toán cần tối thiểu hóa giá trị lớn nhất trong các chi phí này. Cách đặt mục tiêu như vậy phản ánh yêu cầu cân bằng tải giữa các xe tải và drone, vì thời gian hoàn thành toàn bộ nhiệm vụ phụ thuộc vào phương tiện hoàn thành muộn nhất.

Một điểm quan trọng của MM-MT-dLARP là bản chất liên tục của bài toán. Drone không bị giới hạn bởi mạng đường bộ và có thể đi trực tiếp giữa hai điểm bất kỳ trong không gian. Hơn nữa, drone có thể đi vào hoặc rời khỏi một đoạn tuyến tại các điểm nằm bên trong đoạn tuyến, không nhất thiết chỉ tại hai đầu mút. Đặc điểm này giúp tìm được nghiệm tốt hơn so với phương tiện mặt đất, nhưng cũng làm bài toán khó hơn do không gian lựa chọn là liên tục.

Để xử lý bài toán về mặt tính toán, bài báo gốc rời rạc hóa các đoạn tuyến liên tục bằng cách xấp xỉ mỗi đường bằng một chuỗi đa giác gồm hữu hạn các điểm trung gian. Khi đó, drone chỉ được phép đi vào hoặc rời khỏi một đoạn tuyến tại các điểm đã được chọn trước. Bài toán liên tục MM-MT-dLARP được chuyển thành phiên bản rời rạc gọi là \textit{Min--Max Multi--Trip Location Arc Routing Problem} (MM-MT-LARP). Trong phiên bản rời rạc này, các đoạn nhỏ của chuỗi đa giác trở thành các cạnh yêu cầu trong đồ thị, còn các cạnh không yêu cầu biểu diễn việc drone bay trực tiếp giữa hai điểm mà không thực hiện phục vụ.

\section{Tổng quan về bài báo gốc}

Bài báo \textit{The Min--Max Multi--Trip Drone Location Arc Routing Problem} của Corberán, Plana và Sanchis là công trình nền tảng trực tiếp cho đề tài này \cite{corberan2025mmdlarp}. Bài báo đưa ra một lớp bài toán mới kết hợp đồng thời ba yếu tố: định tuyến trên cung, lựa chọn vị trí triển khai, và phối hợp xe tải với drone có khả năng bay nhiều chuyến. Đây là sự kết hợp giữa bài toán Location Arc Routing Problem, bài toán nhiều depot hoặc nhiều điểm vận hành, và bài toán định tuyến drone có giới hạn thời lượng bay.

Đóng góp đầu tiên của bài báo là phát biểu chính thức bài toán MM-MT-dLARP và phiên bản rời rạc MM-MT-LARP. Trong mô hình rời rạc, bài toán được định nghĩa trên đồ thị vô hướng \(G=(V,E)\), trong đó tập cạnh \(E\) gồm các cạnh yêu cầu \(E_R\) và các cạnh không yêu cầu \(E_{NR}\). Các cạnh yêu cầu tương ứng với các đoạn cần drone bay qua để phục vụ, còn các cạnh không yêu cầu biểu diễn các đoạn bay chết, tức là bay từ điểm này sang điểm khác mà không phục vụ. Tập đỉnh bao gồm các điểm thuộc các đoạn tuyến, các điểm triển khai drone, và depot.

Đóng góp thứ hai là xây dựng mô hình quy hoạch nguyên cho MM-MT-LARP. Mô hình sử dụng các biến nhị phân để biểu diễn việc một chuyến bay có đi qua hoặc phục vụ một cạnh hay không, cùng với biến nhị phân biểu diễn việc một điểm triển khai có được mở hay không. Hàm mục tiêu là tối thiểu hóa biến \(t\), đại diện cho chi phí lớn nhất trong các tuyến xe tải--drone. Các ràng buộc chính đảm bảo mỗi cạnh yêu cầu được phục vụ, mỗi chuyến bay là một chu trình liên thông bắt đầu và kết thúc tại điểm triển khai tương ứng, mỗi chuyến bay không vượt quá giới hạn \(L\), và số điểm triển khai được sử dụng không vượt quá số xe tải \(P\).

Đóng góp thứ ba là nghiên cứu đa diện và xây dựng các bất đẳng thức hợp lệ để tăng cường mô hình. Bài báo trình bày nhiều họ bất đẳng thức như ràng buộc liên thông, bất đẳng thức parity, bất đẳng thức (p)-connectivity và bất đẳng thức max-length. Một số họ bất đẳng thức được chứng minh là sinh facet cho một đa diện nới lỏng của bài toán. Trên cơ sở đó, các tác giả phát triển thuật toán branch-and-cut cho phiên bản rời rạc của bài toán.

Đóng góp thứ tư, và cũng là phần liên quan trực tiếp đến đề tài này, là thuật toán matheuristic cho MM-MT-dLARP. Do bài toán có độ khó cao, thuật toán exact chỉ áp dụng hiệu quả cho các thể hiện nhỏ. Vì vậy, bài báo đề xuất một thuật toán heuristic kết hợp giữa pha xây dựng nghiệm ban đầu, các thủ tục tìm kiếm cục bộ, Variable Neighborhood Descent (VND), và một pha chọn điểm trung gian trên các đoạn tuyến.

Cụ thể, trong pha xây dựng nghiệm, thuật toán chọn ngẫu nhiên \(P\) điểm triển khai, sau đó phân công các cạnh yêu cầu cho các điểm triển khai dựa trên xác suất tỉ lệ nghịch với khoảng cách. Với mỗi điểm triển khai, một giant tour được xây dựng theo chiến lược Nearest Neighbor, sau đó được tách thành nhiều chuyến bay khả thi sao cho mỗi chuyến không vượt quá giới hạn \(L\). Sau khi có nghiệm ban đầu, thuật toán áp dụng bốn thủ tục tìm kiếm cục bộ: destroy-and-repair move, intraroute move, \(0\)-to-\(\ell\) exchange, và \(\ell_1\)-to-\(\ell_2\) exchange. Các thủ tục này được tích hợp trong khung Variable Neighborhood Descent để cải thiện nghiệm.

Pha cuối cùng của matheuristic là pha \textit{splitting}, trong đó thuật toán thêm các điểm trung gian vào các đoạn tuyến nhằm khai thác khả năng drone có thể đi vào hoặc rời khỏi một đường tại các điểm nằm bên trong đường. Ban đầu, mỗi đường có thể được biểu diễn bằng một đoạn duy nhất. Sau đó, thuật toán thử thêm trung điểm vào các cạnh, chạy lại VND để kiểm tra xem các điểm trung gian mới có giúp cải thiện nghiệm hay không, rồi giữ lại các điểm hữu ích và loại bỏ các điểm không được sử dụng.

Các thực nghiệm trong bài báo được tiến hành trên tập thể hiện sinh ngẫu nhiên, với tối đa 6 điểm triển khai và 88 đường gốc. Kết quả cho thấy branch-and-cut chỉ phù hợp với các thể hiện nhỏ, trong khi matheuristic có khả năng tìm nghiệm chất lượng tốt trong thời gian hợp lý. Điều này cho thấy đối với MM-MT-dLARP, các phương pháp tìm kiếm cục bộ và metaheuristic đóng vai trò quan trọng khi kích thước bài toán tăng lên.

\section{Các công trình liên quan}

\subsection{Các bài toán định tuyến drone}

Các nghiên cứu về định tuyến drone có thể chia thành hai nhóm chính. Nhóm thứ nhất là các bài toán giao hàng, trong đó drone phục vụ các khách hàng được mô hình hóa dưới dạng các nút. Nhóm này thường liên quan đến các biến thể của Vehicle Routing Problem hoặc Traveling Salesman Problem có drone hỗ trợ. Các tổng quan về drone-aided routing và phối hợp drone với xe tải được trình bày trong \cite{macrina2020droneReview,chung2020droneTruck}.

Nhóm thứ hai là các bài toán kiểm tra hoặc giám sát mạng lưới, trong đó đối tượng phục vụ là các cạnh hoặc cung. Đây là nhóm có liên hệ trực tiếp hơn với MM-MT-dLARP. Campbell và cộng sự đề xuất bài toán Drone Arc Routing Problem, trong đó drone cần phục vụ một tập các đường liên tục và có thể đi vào hoặc rời khỏi đường tại các điểm trung gian \cite{campbell2018darp}. Công trình này đặt nền móng cho việc mô hình hóa bài toán kiểm tra tuyến bằng drone dưới dạng bài toán định tuyến trên cung.

Các nghiên cứu tiếp theo mở rộng hướng này sang bài toán nhiều drone có giới hạn chiều dài chuyến bay. Campbell và cộng sự nghiên cứu bài toán \textit{Length-Constrained (K)-Drones Rural Postman Problem}, trong đó một đội drone cần cùng nhau phục vụ các cạnh yêu cầu dưới ràng buộc độ dài mỗi chuyến bay \cite{campbell2021kdrone}. Sau đó, các tác giả tiếp tục phân tích đa diện và đề xuất thuật toán mới cho bài toán này \cite{campbell2022polyhedral}. Các công trình này có liên hệ chặt chẽ với MM-MT-dLARP vì cùng xử lý đối tượng phục vụ là các cạnh và cùng xét ràng buộc thời lượng bay của drone.

Tuy nhiên, điểm khác biệt của MM-MT-dLARP là bài toán không chỉ thiết kế đường bay cho drone mà còn phải chọn các điểm triển khai cho xe tải. Do đó, bài toán kết hợp thêm quyết định location vào bài toán arc routing, đồng thời sử dụng mục tiêu min--max để cân bằng thời gian hoàn thành giữa các xe tải và drone.

\subsection{Location Arc Routing Problem và bài toán nhiều depot}

Location Arc Routing Problem là lớp bài toán kết hợp giữa quyết định mở cơ sở và quyết định định tuyến trên cung. Trong lớp bài toán này, ta không chỉ cần xác định các tuyến phục vụ các cạnh yêu cầu, mà còn phải quyết định các điểm xuất phát hoặc depot được sử dụng. MM-MT-dLARP có thể được xem là một biến thể đặc biệt của Location Arc Routing Problem, trong đó các cơ sở được mở là các điểm triển khai drone, còn các tuyến tương ứng là các chuyến bay của drone từ các điểm đó.

Bên cạnh yếu tố location, MM-MT-dLARP cũng có liên hệ với các bài toán nhiều depot. Trong các bài toán multi-depot arc routing, các phương tiện xuất phát từ nhiều depot khác nhau và cần phối hợp để phục vụ tập cạnh yêu cầu. Tuy nhiên, trong MM-MT-dLARP, các depot không hoàn toàn cố định từ đầu mà được chọn từ tập điểm triển khai tiềm năng. Ngoài ra, mỗi điểm triển khai gắn với một xe tải và một drone, trong đó drone có thể thực hiện nhiều chuyến bay. Điều này làm tăng độ phức tạp so với các bài toán nhiều depot cổ điển.

\subsection{Các phương pháp tìm kiếm cục bộ cho bài toán định tuyến}

Do các bài toán định tuyến thường có độ phức tạp tổ hợp cao, các phương pháp tìm kiếm cục bộ và metaheuristic được sử dụng rộng rãi để tìm nghiệm tốt trong thời gian hợp lý. Tìm kiếm cục bộ bắt đầu từ một nghiệm khả thi và lặp lại việc di chuyển sang nghiệm lân cận tốt hơn theo một tập phép biến đổi được định nghĩa trước. Chất lượng của thuật toán phụ thuộc nhiều vào cách mã hóa nghiệm, cách định nghĩa lân cận, chiến lược chọn nước đi và cơ chế thoát khỏi cực trị địa phương.

Variable Neighborhood Search (VNS) là một metaheuristic dựa trên ý tưởng thay đổi có hệ thống các cấu trúc lân cận trong quá trình tìm kiếm \cite{mladenovic1997vns}. Thay vì chỉ sử dụng một loại lân cận cố định, VNS khai thác nhiều lân cận khác nhau để tăng khả năng thoát khỏi cực trị địa phương. Một biến thể quan trọng của VNS là Variable Neighborhood Descent (VND), trong đó thuật toán lần lượt áp dụng các thủ tục tìm kiếm cục bộ trên các lân cận khác nhau cho đến khi không còn cải thiện. Các nguyên lý và ứng dụng của VNS được trình bày trong \cite{hansen2001vns,hansen2010vns}.

Large Neighborhood Search (LNS) là một hướng tiếp cận khác, trong đó thuật toán phá hủy một phần nghiệm hiện tại rồi xây dựng lại phần bị phá hủy để tạo nghiệm mới \cite{shaw1998lns}. Cơ chế destroy-and-repair giúp LNS khám phá các lân cận lớn hơn nhiều so với các phép hoán đổi nhỏ truyền thống. Adaptive Large Neighborhood Search (ALNS) mở rộng LNS bằng cách sử dụng nhiều toán tử destroy và repair khác nhau, đồng thời điều chỉnh xác suất chọn toán tử dựa trên hiệu quả của chúng trong quá trình tìm kiếm \cite{ropke2006alns,pisinger2007vrp}.

Trong bài báo gốc về MM-MT-dLARP, các tác giả đã sử dụng một khung VND gồm bốn thủ tục cải thiện nghiệm. Trong đó, destroy-and-repair move có bản chất gần với LNS vì nó xóa một số cạnh yêu cầu khỏi nghiệm hiện tại rồi thử chèn lại vào các vị trí tốt hơn. Tuy nhiên, thuật toán trong bài báo chưa phải là một VNS hoặc LNS đầy đủ theo nghĩa cổ điển. Đây chính là cơ sở để đề tài này tập trung nghiên cứu và phát triển các giải thuật tìm kiếm cục bộ, VNS và LNS cho bài toán MM-MT-dLARP.

\section{Định hướng của đồ án}

Dựa trên bài báo gốc, đồ án này tập trung vào khía cạnh thuật toán tìm kiếm cục bộ cho bài toán MM-MT-dLARP. Thay vì đi sâu vào mô hình quy hoạch nguyên và branch-and-cut, đồ án hướng đến việc phân tích cách biểu diễn nghiệm, các cấu trúc lân cận, và khả năng xây dựng các khung metaheuristic như VNS và LNS.

Cụ thể, đồ án xem xét nghiệm của bài toán dưới dạng tập các điểm triển khai được chọn và tập các chuyến bay của drone tại mỗi điểm triển khai. Một nghiệm tốt cần thỏa mãn đồng thời ba yếu tố: mỗi cạnh yêu cầu được phục vụ, mỗi chuyến bay không vượt quá giới hạn \(L\), và giá trị makespan được giảm thiểu. Từ biểu diễn này, có thể xây dựng nhiều nhóm toán tử lân cận như di chuyển một cạnh yêu cầu giữa hai chuyến bay, đổi chỗ hai chuỗi cạnh giữa hai chuyến bay, tối ưu lại thứ tự phục vụ trong một chuyến bay, hoặc phá hủy và tái xây dựng một phần nghiệm tại điểm triển khai đang là nút thắt.

Đối với hướng VNS, đồ án có thể định nghĩa nhiều cấu trúc lân cận \(N_1, N_2, \ldots, N_k\), ví dụ như relocate, swap, sequence exchange và destroy-and-repair. Quá trình tìm kiếm sẽ thay đổi lân cận một cách có hệ thống, kết hợp bước shaking để tạo nhiễu và bước local search để cải thiện nghiệm. Đối với hướng LNS, đồ án có thể tập trung vào các chiến lược destroy-and-repair chuyên biệt cho MM-MT-dLARP, chẳng hạn ưu tiên phá hủy tại điểm triển khai có makespan lớn nhất, xóa các cạnh có chi phí chèn cao, hoặc tái chèn các cạnh theo tiêu chí giảm bottleneck.

Như vậy, chương này đã trình bày tổng quan về bối cảnh ứng dụng, mô tả khái quát bài toán MM-MT-dLARP, tóm tắt nội dung chính của bài báo nền tảng và các hướng nghiên cứu liên quan. Các chương tiếp theo sẽ đi sâu hơn vào mô hình bài toán, cách mã hóa nghiệm, các toán tử lân cận và thiết kế thuật toán tìm kiếm cục bộ cho bài toán này.
